\section*{Acknowledgement of AI use}
Artificial Intelligence (AI) tools were used in the structuring and writing of this project plan and time plan. The idea and outline of the project was developed by the student and supervisor without the use of AI tools. Figures were created by the student without the use of AI tools. The student and supervisor ensures physical accuracy and validity of the text.

\section{Project Plan}
Quantum measurement is a fundamental aspect of quantum mechanics and is particularly important for the readout and control of solid-state quantum systems. This project will theoretically investigate quantum measurements of quantum-dot(s) systems coupled to microwave resonators, building on previous work within the research group (Nord, 2026) and in close contact with experimental colleagues in the group of Ville Maisi, where related set-ups have been analyzed. In the experiment, information about the quantum-dot(s) system is obtained by probing a microwave resonator and measuring its reflection coefficient.

The project will initially consider a single quantum dot as a relatively simple system in which the connection between the state of the device and the microwave response can be studied. The main focus will then be shifted towards a double quantum dot (DQD), where the relevant degrees of freedom can form an effective two-level system, or qubit. The aim is to understand how a quantum state can be measured through its coupling to a hybrid microwave system and how the measurement itself affects the quantum system.

\begin{figure}[h]
    \centering
    \includegraphics{system.pdf}
    \caption{The main system that will be considered. The left part shows the hybrid microwave system that will act as a detector. The right part shows a two level system (which will be modelled as a double quantum dot) and the capacitive coupling between the two systems.}
\end{figure}

A central theoretical framework will be input-output theory, which relates the internal dynamics of a microwave resonator to the experimentally accessible input and output fields. For a resonator coupled to an input field, the output field can, in a simplified description, be written as

$$
b_{\mathrm{out}}(t)=b_{\mathrm{in}}(t)+\sqrt{\kappa_{\mathrm{c}}}\,a(t),
$$

where \(a\) is the cavity field and \(\kappa_{\mathrm{c}}\) is the coupling rate to the external transmission line. The measured reflection coefficient can consequently be related to the state of the quantum-dot system through its influence on the resonator response.

The coupling between the DQD and the resonator can be described by the Jaynes-Cummings Hamiltonian, or more generally the full quantum Rabi model. Both resonant and dispersive coupling regimes will be considered and are of interest for the experimental group.

An important part of the project will be to understand the measurement process qualitatively, complementing mathematical descriptions that have already been developed within the group. Particular attention will be given to measurement backaction, relaxation and decoherence, and to the conditions under which the measurement can be considered quantum non-demolition (QND). The detector sensitivity and the amount of information obtainable from the measurement will also be investigated. Fisher information will be considered as a possible measure of the information contained in the measured microwave signal about a system parameter $\theta$.


A further objective is to investigate the approximations used in the theoretical description and determine under which physical conditions they are valid. 